\documentclass[12pt]{article}
\usepackage[utf8]{inputenc}
\usepackage[T1]{fontenc}
\usepackage[margin=1in]{geometry}
\usepackage{setspace}
\usepackage{xcolor}
\usepackage[colorlinks=true, urlcolor=blue!60!black, linkcolor=blue!60!black]{hyperref}
\setlength{\parindent}{0em}
\setlength{\parskip}{0.6em}
\renewcommand{\baselinestretch}{1.0}
\begin{document}
\begin{center}
{\LARGE \textit{Chaplinsky v. New Hampshire}, 315 U.S. 568 (1942)\par}
\vspace{0.5em}
{\large Case Brief --- Two-Page Edition\par}
\end{center}

\section{Memory Jogger}
``Fighting words,'' face-to-face epithets likely to provoke an average person to violence, are unprotected; a Jehovah's Witness's conviction for calling a marshal a ``damned racketeer'' and ``damned Fascist'' stands (\href{https://supreme.justia.com/cases/federal/us/315/568/}{Justia}).

\section{Detailed Case Facts}
Walter Chaplinsky, a Jehovah's Witness, denounced organized religion as a ``racket'' on a Rochester, New Hampshire, street, and a crowd grew restless. While being escorted toward the police station, he met City Marshal Bowering and called him ``a God damned racketeer'' and ``a damned Fascist.''

\section{Procedural History}
He was convicted under a state law barring ``offensive, derisive or annoying'' words to persons in public. The New Hampshire Supreme Court affirmed, construing the law to reach only face-to-face fighting words. The U.S. Supreme Court affirmed in 1942.

\section{Judicial Votes}
9--0. Justice Murphy wrote for a unanimous Court; there were no separate opinions.

\section{Holding}
Narrowly limited classes of speech, including fighting words that ``by their very utterance inflict injury or tend to incite an immediate breach of the peace,'' are unprotected. As narrowly construed, the statute is neither vague nor an abridgment of free speech.

\section{Analysis of Opinions}
\textbf{Majority:} Free speech is ``not absolute at all times and under all circumstances.'' The lewd, obscene, profane, libelous, and fighting words are ``no essential part of any exposition of ideas,'' and their slight value is outweighed by order and morality. The state court's objective test, what ``men of common intelligence'' would see as likely to provoke a fight, gave fair notice. \textbf{Later:} \textit{Cohen} and \textit{Gooding} narrowed the category to face-to-face epithets risking immediate violence, \textit{R.A.V.} barred viewpoint discrimination within it, and \textit{Stevens} treated it as a closed historical category.

\section{Examples: Future Applications}
\textbf{Same-Side:} personal slurs and a dare to fight delivered inches from a neighbor's face are unprotected. \textbf{Opposite-Side:} insults shouted at police through a bullhorn or posted online are protected. \textbf{Fence-Sitter:} Chaplinsky's own words spoken to an arresting officer today.

\section{Critique}
Chaplinsky's words were political criticism of an official who failed to protect a religious minority, speech modern doctrine protects. The ``average addressee'' assumption invites selective enforcement, and the ``inflict injury'' prong conflicts with \textit{Snyder v. Phelps}.

\section{Key Quotations}
\begin{quote}
``These include the lewd and obscene, the profane, the libelous, and the insulting or `fighting' words---those which by their very utterance inflict injury or tend to incite an immediate breach of the peace. It has been well observed that such utterances are no essential part of any exposition of ideas, and are of such slight social value as a step to truth that any benefit that may be derived from them is clearly outweighed by the social interest in order and morality.''
\end{quote}
\noindent --- Justice Murphy, for the Court, 315 U.S. at 572.

\end{document}
